\section{Goldilocks Field Arithmetic}
\label{sec:goldilocks}

The polynomial evaluation methods of Section~\ref{sec:poly-eval} ultimately
depend on the cost of field multiplication.  We therefore separate the
choice of polynomial representation from the choice of modular
multiplication algorithm.

The experiments use the Goldilocks prime
\[
p=2^{64}-2^{32}+1.
\]
A field element is represented by one machine word.  Two multiplication
strategies are considered.  The first exploits the special form of the
modulus.  The second uses Montgomery representation with radix
\[
R=2^{64}.
\]
Both strategies are retained in the experiments so that the effect of the
modular multiplication algorithm can be measured independently of the
polynomial representation.

\subsection{Reduction Specialized to the Goldilocks Modulus}
\label{sec:goldilocks-special}

Let
\[
z=ab
\]
be the \(128\)-bit product of two canonical residues
\[
0\le a,b<p.
\]
Write
\[
z=z_{\mathrm{lo}}+2^{64}z_{\mathrm{hi}},
\]
where
\[
0\le z_{\mathrm{lo}},z_{\mathrm{hi}}<2^{64}.
\]

For the Goldilocks prime,
\[
2^{64}\equiv 2^{32}-1\pmod p.
\]
Define
\[
\varepsilon=2^{32}-1.
\]
Then
\[
z
\equiv
z_{\mathrm{lo}}+\varepsilon z_{\mathrm{hi}}
\pmod p.
\]

A direct application of this congruence already reduces the \(128\)-bit
product to a value close to one machine word.  A more useful form is
obtained by splitting the high half as
\[
z_{\mathrm{hi}}
=
h_0+2^{32}h_1,
\qquad
0\le h_0,h_1<2^{32}.
\]
Since
\[
2^{32}\varepsilon
=
2^{64}-2^{32}
\equiv -1
\pmod p,
\]
we obtain
\[
\varepsilon z_{\mathrm{hi}}
\equiv
h_0\varepsilon-h_1
\pmod p.
\]
Therefore
\[
z
\equiv
z_{\mathrm{lo}}
-h_1
+h_0(2^{32}-1)
\pmod p.
\]

This expression involves only the \(128\)-bit product \(ab\), extraction of
the two \(32\)-bit halves of \(z_{\mathrm{hi}}\), shifts, additions, and
subtractions, followed by a small number of corrections to return to the
canonical interval.

The implementation used in the benchmark follows this reduction rather than
a generic integer remainder instruction.

\paragraph{Algorithm 1. Goldilocks multiplication.}

\[
\begin{array}{ll}
\textbf{Input:} & a,b\in[0,p),\\
\textbf{Output:} & ab\bmod p.
\end{array}
\]

\[
\begin{array}{rcl}
z &\leftarrow& a\cdot b \qquad\text{(full \(128\)-bit product)},\\
z_{\mathrm{lo}} &\leftarrow& z\bmod 2^{64},\\
z_{\mathrm{hi}} &\leftarrow& \lfloor z/2^{64}\rfloor,\\
h_0 &\leftarrow& z_{\mathrm{hi}}\bmod 2^{32},\\
h_1 &\leftarrow& \lfloor z_{\mathrm{hi}}/2^{32}\rfloor,\\
u &\leftarrow& z_{\mathrm{lo}}-h_1+h_0(2^{32}-1),\\
u &\leftarrow& \text{canonical reduction of }u\text{ modulo }p,\\
\textbf{return}&&u.
\end{array}
\]

The exact instruction sequence generated from this formula depends on the
compiler and target architecture.  For this reason the multiplication is
measured separately before its effect on the complete polynomial evaluator
is considered.

\subsection{Montgomery Representation}
\label{sec:montgomery}

Montgomery multiplication replaces division by \(p\) with arithmetic modulo
a power of two.  Let
\[
R=2^{64}.
\]
Because \(p\) is odd,
\[
\gcd(p,R)=1,
\]
and \(R\) is invertible modulo \(p\).

A canonical field element \(a\) is represented in Montgomery form by
\[
\widetilde{a}=aR\bmod p.
\]
Multiplication of two Montgomery residues produces
\[
\widetilde{a}\widetilde{b}
=
abR^2
\]
and is followed by a Montgomery reduction, which removes one factor of
\(R\):
\[
\operatorname{REDC}
\left(
\widetilde{a}\widetilde{b}
\right)
=
abR\bmod p.
\]

For the modulus considered here, the constants are
\[
R\bmod p
=
2^{32}-1
=
\mathtt{0x00000000ffffffff},
\]
\[
n'=-p^{-1}\bmod R
=
\mathtt{0xfffffffeffffffff},
\]
and
\[
R^2\bmod p
=
\mathtt{0xfffffffe00000001}.
\]

The value \(R^2\bmod p\) is used to convert a canonical residue into
Montgomery form with one Montgomery multiplication.

\subsection{Montgomery Reduction}
\label{sec:montgomery-redc}

Let
\[
0\le T<pR.
\]
The Montgomery reduction computes
\[
TR^{-1}\bmod p
\]
without an integer division by \(p\).

Using
\[
n'=-p^{-1}\bmod R,
\]
define
\[
m=(T\bmod R)n'\bmod R.
\]
By construction,
\[
T+mp
\equiv 0\pmod R,
\]
so the quotient
\[
u=\frac{T+mp}{R}
\]
is an integer.  The result satisfies
\[
u\equiv TR^{-1}\pmod p.
\]
A final conditional subtraction returns the canonical representative.

\paragraph{Algorithm 2. Montgomery reduction.}

\[
\begin{array}{ll}
\textbf{Input:} & T,\quad 0\le T<pR,\\
\textbf{Output:} & TR^{-1}\bmod p.
\end{array}
\]

\[
\begin{array}{rcl}
m &\leftarrow& (T\bmod R)n'\bmod R,\\
u &\leftarrow& (T+mp)/R,\\
\textbf{if }u\ge p && u\leftarrow u-p,\\
\textbf{return}&&u.
\end{array}
\]

For a \(64\)-bit radix, the sum
\[
T+mp
\]
may require \(129\) bits.  The implementation therefore records the carry
from the \(128\)-bit addition explicitly before extracting the high word.

\subsection{Montgomery Multiplication}
\label{sec:montgomery-mul}

If \(\widetilde{a}\) and \(\widetilde{b}\) are already stored in Montgomery
form, their product is computed as
\[
\operatorname{MontMul}
(\widetilde{a},\widetilde{b})
=
\operatorname{REDC}
(\widetilde{a}\widetilde{b}).
\]

\paragraph{Algorithm 3. Montgomery multiplication.}

\[
\begin{array}{ll}
\textbf{Input:} & \widetilde{a}=aR\bmod p,\quad
                  \widetilde{b}=bR\bmod p,\\
\textbf{Output:} & abR\bmod p.
\end{array}
\]

\[
\begin{array}{rcl}
T &\leftarrow& \widetilde{a}\widetilde{b},\\
u &\leftarrow& \operatorname{REDC}(T),\\
\textbf{return}&&u.
\end{array}
\]

Conversion from canonical to Montgomery form is performed as
\[
\widetilde{a}
=
\operatorname{REDC}
\left(
a(R^2\bmod p)
\right),
\]
while conversion back to the canonical representation is
\[
a
=
\operatorname{REDC}(\widetilde{a}).
\]

In the multiplication benchmark, operands are kept in their native
representation throughout the timed region.  Thus the specialized
Goldilocks implementation receives canonical residues, while the
Montgomery implementation receives Montgomery residues.  Conversion is
measured separately when it is relevant to an application and is not
charged once per multiplication.

\subsection{Kernel Fold Under the Two Arithmetic Models}
\label{sec:fold-arithmetic}

The Boolean-kernel reduction uses
\[
F_t(a,b)=t(a+b)+b.
\]
With canonical Goldilocks residues, the implementation can combine the
final addition by \(b\) with the reduction of the \(128\)-bit product.
This gives a specialized field operation rather than a literal sequence of
one modular addition, one modular multiplication, and a second modular
addition.

In Montgomery form the same mathematical operation is
\[
\widetilde{F}_t(\widetilde{a},\widetilde{b})
=
\operatorname{MontMul}
\left(
\widetilde{t},
\widetilde{a}+\widetilde{b}
\right)
+
\widetilde{b}.
\]
All operands remain in Montgomery form, so no conversion is needed between
levels of the reduction tree.

Table~\ref{tab:field-mul-structure} summarizes the two multiplication
strategies before measurements.

\begin{table}[ht]
\centering
\caption{Structure of the two Goldilocks multiplication strategies.}
\label{tab:field-mul-structure}
\begin{tabular}{lll}
\toprule
Method & Representation & Main reduction operation\\
\midrule
Goldilocks special &
canonical residue &
use \(2^{64}\equiv2^{32}-1\pmod p\)\\

Montgomery &
\(aR\bmod p\) &
\((T+mp)/2^{64}\) with \(m=T_0n'\)\\
\bottomrule
\end{tabular}
\end{table}

At this point there is no architectural conclusion to draw from the
operation formulas alone.  The specialized reduction uses the algebraic
form of the Goldilocks modulus; Montgomery reduction introduces additional
word multiplications but has a regular reduction structure.  The generated
instruction sequences and the complete evaluation times are therefore
examined experimentally in the following sections.
